\documentclass[11pt]{article}
\usepackage[margin=1in]{geometry}
\usepackage{amsmath,amssymb,amsthm}
\usepackage{graphicx}
\usepackage{microtype}
\usepackage[colorlinks=true,linkcolor=blue,citecolor=blue,urlcolor=blue]{hyperref}

\newtheorem{theorem}{Theorem}
\newtheorem{lemma}{Lemma}
\newtheorem{remark}{Remark}
\newcommand{\U}{\mathcal U}
\newcommand{\subsetsim}{\mathrel{\hookrightarrow}}

\title{A counterexample to Kruglov necessity for embeddings of
generalized Rosenthal spaces}
\author{Candidate counterexample for arXiv:2006.00936}
\date{11 August 2026}

\begin{document}
\maketitle

\begin{center}
\fbox{\parbox{0.91\linewidth}{\textbf{Status.} Candidate counterexample,
likely valid, requiring human review. In the unrestricted formulation printed
in the source, the Kruglov property is not necessary: take
$E=L^\infty[0,1]$.}}
\end{center}

\section{Source problem}

For a rearrangement-invariant Banach function space $E$ on $[0,1]$,
Astashkin and Curbera define a rearrangement-invariant space $Z_E$ on
$(0,\infty)$ and then define the generalized Rosenthal space $\U_E$ as the
closed linear span in $Z_E$ of a countable sequence of characteristic
functions $\{\chi_{A_n}\}_{n=1}^\infty$ satisfying their conditions (2).
Their notation $Y\subsetsim X$ means that there is an arbitrary isomorphic
embedding from $Y$ into $X$.

The source says on PDF page 4:
\begin{quote}
``In the concluding part of the paper, in Theorem 6, we prove a partial
result related to the problem if the Kruglov property of a r.i. space $E$ is
a necessary condition for the existence of an isomorphic embedding
$T:\U_E\to E$.''
\end{quote}
The stated partial result assumes that images of suitable canonical block
vectors are independent symmetrically distributed random variables satisfying
an additional uniform lower-tail condition \cite[Theorem~6]{AC2022}.

\begin{figure}[ht]
\centering
\includegraphics[width=\linewidth]{figures/open_problem_crop.png}
\caption{The problem and the scope of the source's partial result on PDF
page 4 of \cite{AC2022}.}
\end{figure}

\section{Proof intuition}

The issue is that an arbitrary Banach-space embedding need not retain any of
the probabilistic structure of the canonical basis. The space $\U_E$ is
always separable, simply because it is generated by a countable sequence.
At the opposite extreme, $L^\infty[0,1]$ is a universal host for separable
Banach spaces, so it automatically contains a copy of its own $\U_E$.

But $L^\infty$ is too close to the endpoint to have the Kruglov property:
the Kruglov transform of the constant one function is an unbounded Poisson
random variable. These two elementary observations give the counterexample.

\section{Two elementary embedding facts}

\begin{lemma}\label{lem:separable}
For every rearrangement-invariant space $E$ in the source construction,
$\U_E$ is separable.
\end{lemma}

\begin{proof}
By definition, $\U_E$ is the closed linear span of the countable sequence
$\{\chi_{A_n}\}_{n=1}^\infty$. Finite linear combinations using a countable
dense subfield of the scalars form a countable dense subset of this closed
span.
\end{proof}

\begin{lemma}\label{lem:universal}
Every separable Banach space embeds linearly isometrically into
$L^\infty[0,1]$.
\end{lemma}

\begin{proof}
Let $X$ be separable. If $X=\{0\}$ there is nothing to prove. Choose a dense
sequence $(x_n)$ in the unit sphere of $X$. By Hahn--Banach, choose
$x_n^*\in X^*$ with
\[
                    \|x_n^*\|=1,\qquad x_n^*(x_n)=1.
\]
In the complex case the phase may be chosen so that the second equality is
literal. Define
\[
 J:X\longrightarrow\ell_\infty,\qquad
 Jx=(x_n^*(x))_{n=1}^\infty.
\]
This map is linear and contractive. If $\|x\|=1$ and $x_{n_k}\to x$, then
\[
 |x_{n_k}^*(x)|
 \ge |x_{n_k}^*(x_{n_k})|-\|x-x_{n_k}\|
 =1-\|x-x_{n_k}\|\longrightarrow1.
\]
Therefore $\|Jx\|_\infty=1$, and homogeneity shows that $J$ is an isometry.

Now choose a measurable partition $(B_n)_{n=1}^\infty$ of $[0,1]$ with
$m(B_n)>0$ for every $n$, and define
\[
 S:\ell_\infty\longrightarrow L^\infty[0,1],\qquad
 S((a_n))=\sum_{n=1}^\infty a_n\chi_{B_n}.
\]
Every coordinate value occurs on a set of positive measure, so
$\|S((a_n))\|_{L^\infty}=\sup_n|a_n|$. Hence $S$ and $S\circ J$ are linear
isometries.
\end{proof}

\section{The counterexample}

Recall that $E$ has the \emph{Kruglov property} if $f\in E$ implies
$\pi(f)\in E$, where
\[
                    \pi(f)=\sum_{i=1}^{N}f_i,
\]
the $f_i$ are independent copies of $f$, and $N$ is an independent Poisson
random variable with parameter one.

\begin{theorem}\label{thm:main}
The Kruglov property is not necessary for the existence of an isomorphic
embedding $\U_E\to E$ in the unrestricted formulation of
\cite{AC2022}. More precisely, the rearrangement-invariant space
\[
                           E=L^\infty[0,1]
\]
has the Fatou property and satisfies $\U_E\subsetsim E$, but it does not have
the Kruglov property.
\end{theorem}

\begin{proof}
Lemma~\ref{lem:separable} says that $\U_{L^\infty}$ is separable. Applying
Lemma~\ref{lem:universal} with $X=\U_{L^\infty}$ gives a linear isometric
embedding
\[
                    \U_{L^\infty}\longrightarrow L^\infty[0,1].
\]
Thus the required isomorphic embedding exists. The familiar monotone-limit
property of the essential-supremum norm shows that $L^\infty$ has the Fatou
property.

To disprove the Kruglov property, take $f=\mathbf1$. Every independent copy
of $f$ is again the constant one function, and hence
\[
                         \pi(\mathbf1)=N.
\]
For every integer $m\ge0$,
\[
                         \mathbb P(N=m)=\frac{e^{-1}}{m!}>0.
\]
Consequently $N$ is essentially unbounded. Thus
$\mathbf1\in L^\infty[0,1]$ while
$\pi(\mathbf1)\notin L^\infty[0,1]$, and $L^\infty$ has no Kruglov
property.
\end{proof}

\begin{remark}[Scope of the counterexample]
The embedding above is deliberately arbitrary; it need not send canonical
vectors to independent random variables or satisfy the lower-tail hypothesis
of \cite[Theorem~6]{AC2022}. There is therefore no conflict with that partial
theorem. If the intended question silently restricts $E$ to separable or
order-continuous r.i. spaces, or restricts $T$ to preserve canonical
probabilistic structure, that strengthened question remains outside the
present result. Those restrictions do not occur in the printed introduction
statement.
\end{remark}

\section{Verification and bounded novelty check}

The proof is noncomputational. Its only ingredients are the source definition
of $\U_E$, a countable Hahn--Banach norming family, a positive-measure
partition of $[0,1]$, and the elementary Poisson law. The accompanying
\texttt{verification.md} checks each step and the exact scope.

On 11 August 2026, the local run indexes and downloaded arXiv corpus were
searched by arXiv id and the terms ``Kruglov property'', ``$\U_E$'',
``isomorphic embedding'', and ``necessary''. External exact and close
searches used the same phrases, the paper title and authors, and the journal
DOI. They found the source and earlier sufficient-condition literature but no
explicit answer or this $L^\infty$ counterexample. OpenAlex listed no citing
works for the DOI at the time of checking; MaRDI listed five related citing
titles, none of which surfaced an answer in the bounded search. Novelty
confidence is moderate and requires specialist confirmation.

\paragraph{Human-review recommendation.}
Confirm the unrestricted meaning of $\subsetsim$ and decide whether the
authors intended an unstated separability or structure-preserving condition.
No mathematical gap is known for the literal formulation.

\begin{thebibliography}{9}
\bibitem{AC2022}
S. V. Astashkin and G. P. Curbera,
\emph{Rosenthal's space revisited},
Studia Math. \textbf{262} (2022), no.~2, 197--224,
\href{https://doi.org/10.4064/sm201011-4-1}{doi:10.4064/sm201011-4-1};
\href{https://arxiv.org/abs/2006.00936}{arXiv:2006.00936}.
\end{thebibliography}

\end{document}
